\documentclass[12pt,a4paper]{article}

\usepackage[utf8]{inputenc}
\usepackage[T1]{fontenc}
\usepackage[spanish,es-nodecimaldot]{babel}
\usepackage[a4paper,margin=2.5cm]{geometry}
\usepackage{lmodern}
\usepackage{xcolor}
\usepackage{listings}
\usepackage{booktabs}
\usepackage{array}
\usepackage{hyperref}
\usepackage{enumitem}

\hypersetup{
  colorlinks=true,
  linkcolor=blue!50!black,
  urlcolor=blue!60!black,
  pdftitle={Reporte de práctica DevOps},
  pdfauthor={Estudiante}
}

\definecolor{codebg}{RGB}{246,247,249}
\definecolor{codeframe}{RGB}{210,215,222}
\definecolor{codeblue}{RGB}{35,75,145}
\definecolor{codegreen}{RGB}{35,115,75}

\lstdefinestyle{codigo}{
  backgroundcolor=\color{codebg},
  frame=single,
  rulecolor=\color{codeframe},
  basicstyle=\ttfamily\small,
  keywordstyle=\color{codeblue},
  commentstyle=\color{codegreen},
  breaklines=true,
  columns=fullflexible,
  keepspaces=true,
  showstringspaces=false,
  tabsize=2
}

\setlength{\parskip}{0.5em}
\setlength{\parindent}{0pt}

\title{Reporte de práctica: API REST y pruebas automatizadas}
\author{\textbf{Estudiante:} Nombre del estudiante\\
        \textbf{Asignatura:} Gestión y DevOps\\
        \textbf{Docente:} Nombre del docente}
\date{\today}

\begin{document}

\maketitle

\begin{abstract}
En esta práctica se implementó y verificó una API REST para administrar elementos. El servicio está desarrollado con Node.js y Express, utiliza SQLite para la persistencia de datos y se empaqueta mediante Docker. Se configuraron pruebas automatizadas con Jest y Supertest para validar operaciones HTTP, validaciones de entrada y respuestas ante recursos inexistentes. La ejecución final de la suite completó satisfactoriamente las diez pruebas.
\end{abstract}

\section{Objetivos}
\begin{itemize}[leftmargin=*]
  \item Implementar una API HTTP para crear, consultar, actualizar y eliminar elementos.
  \item Persistir los registros en una base de datos SQLite.
  \item Verificar los casos exitosos y los errores esperados mediante pruebas automatizadas.
  \item Describir la configuración de ejecución local y en contenedor.
\end{itemize}

\section{Tecnologías utilizadas}
\begin{itemize}[leftmargin=*]
  \item \textbf{Node.js 18}: entorno de ejecución de JavaScript.
  \item \textbf{Express 5}: definición de rutas HTTP.
  \item \textbf{SQLite}: almacenamiento de elementos.
  \item \textbf{Jest y Supertest}: ejecución de pruebas y solicitudes HTTP.
  \item \textbf{Docker}: construcción y ejecución del servicio en un contenedor basado en Alpine.
\end{itemize}

\section{Descripción de la solución}
La API administra recursos \texttt{items}, cada uno con un identificador, un nombre y una descripción. La base de datos se crea o inicializa al iniciar el servidor. El servidor HTTP escucha en el puerto \texttt{80}; además, el proyecto contiene un servidor TCP en el puerto \texttt{6061}.

Las rutas HTTP principales utilizadas en la práctica son:
\begin{center}
\begin{tabular}{@{}ll@{}}
\toprule
\textbf{Método y ruta} & \textbf{Operación} \\
\midrule
\texttt{GET /items} & Consultar todos los elementos \\
\texttt{GET /items/:id} & Consultar un elemento por identificador \\
\texttt{POST /items} & Crear un elemento \\
\texttt{PUT /items/:id} & Actualizar uno o más campos \\
\texttt{DELETE /items/:id} & Eliminar un elemento \\
\bottomrule
\end{tabular}
\end{center}

Las solicitudes de creación requieren \texttt{name} y \texttt{description}. La actualización permite enviar uno de esos campos, pero rechaza un cuerpo vacío o valores de texto vacíos. Si se consulta, actualiza o elimina un identificador inexistente, la API responde con estado HTTP \texttt{404}.

\section{Pruebas automatizadas}
Las pruebas se encuentran en \texttt{api.test.js}. Supertest envía solicitudes al servidor local, y Jest comprueba los códigos de respuesta y la estructura básica de los datos. La URL de destino predeterminada es \url{http://localhost:80}; puede cambiarse mediante la variable de entorno \texttt{API\_URL}.

\begin{center}
\begin{tabular}{@{}clp{8.2cm}c@{}}
\toprule
\textbf{N.º} & \textbf{Método} & \textbf{Escenario} & \textbf{Resultado} \\
\midrule
1 & GET & Consulta de la colección \texttt{/items}. & Aprobada \\
2 & POST & Creación de un elemento válido. & Aprobada \\
3 & POST & Rechazo de un cuerpo vacío. & Aprobada \\
4 & GET & Consulta del elemento recién creado. & Aprobada \\
5 & GET & Respuesta \texttt{404} para un identificador inexistente. & Aprobada \\
6 & PUT & Actualización del elemento. & Aprobada \\
7 & PUT & Rechazo de un nombre vacío. & Aprobada \\
8 & DELETE & Eliminación del elemento creado. & Aprobada \\
9 & DELETE & Respuesta \texttt{404} al eliminarlo nuevamente. & Aprobada \\
10 & GET & Respuesta \texttt{404} para una ruta inexistente. & Aprobada \\
\bottomrule
\end{tabular}
\end{center}

\textbf{Resultado de ejecución:} 1 suite aprobada, 10 pruebas aprobadas y 0 fallidas.

\subsection{Ejecución local}
En una terminal, se inicia el servidor:
\begin{lstlisting}[style=codigo]
npm start
\end{lstlisting}

En otra terminal, desde la raíz del proyecto, se ejecutan las pruebas:
\begin{lstlisting}[style=codigo]
npm test
\end{lstlisting}

El servidor debe permanecer activo mientras se ejecuta la suite.

\section{Contenedorización}
El \texttt{Dockerfile} usa la imagen \texttt{node:18-alpine}, copia los manifiestos de npm, instala dependencias, copia el proyecto y expone los puertos \texttt{80} y \texttt{6061}. La imagen inicia el servicio con \texttt{node server.js}.

Comandos para construir y ejecutar la imagen:
\begin{lstlisting}[style=codigo]
docker build -t api-items .
docker run --rm -p 80:80 -p 6061:6061 api-items
\end{lstlisting}

\textbf{Nota:} el Dockerfile no configura un volumen para SQLite. Por lo tanto, para conservar la base de datos al eliminar el contenedor se debe configurar almacenamiento persistente.

\section{Conclusiones}
La práctica permitió integrar una API REST con persistencia SQLite y pruebas automatizadas. La suite valida operaciones CRUD, reglas básicas de entrada y respuestas de error. La ejecución satisfactoria de las diez pruebas confirma que los comportamientos cubiertos están implementados conforme a las expectativas definidas en el proyecto. Docker proporciona una forma reproducible de empaquetar el servicio, mientras que la persistencia requiere configurar un volumen cuando los datos deban sobrevivir al ciclo de vida del contenedor.

\end{document}
